%% Chapter 8 — Speed-Enabled Phase Transitions in Reasoning Capability

\chapter{Speed-Enabled Phase Transitions in Reasoning Capability}
\label{chap:speed}

\epigraph{%
  ``Speed is the ultimate weapon. All other factors being equal, the fastest will win.''
}{--- Boyd, OODA loop doctrine; adapted for symbolic reasoning ensembles}

The preceding chapters established what reasoning systems must satisfy: categorical
well-formedness (Chapter~\ref{chap:breed}), oracle-verifiable output (Chapter~\ref{chap:oracle}),
OCEL process conformance (Chapter~\ref{chap:ocel}), BVC admission control
(Chapter~\ref{chap:bvc}), and the fourth causal constraint (Chapter~\ref{chap:fourth-constraint}).
What none of those chapters addressed is the \emph{quantitative} question: how does
execution latency determine \emph{which} of those properties become computationally
accessible within a fixed decision budget?

This chapter proves that below a critical latency threshold $\tau^*$, three qualitatively
distinct reasoning capabilities emerge as phase transitions: Bayesian sensitivity analysis
across a full breed ensemble, Pareto-frontier multi-objective planning, and adversarial
robustness auditing.  These are not merely faster versions of existing capabilities; they
are capabilities that are \emph{logically impossible} when execution latency is above $\tau^*$.

%% -----------------------------------------------------------------------
\section{The Ensemble Size Function}
\label{sec:ensemble-size}
%% -----------------------------------------------------------------------

Fix a decision context.  Let $T_D$ be the available decision budget (wall-clock time
allocated to the reasoning system before an answer must be produced) and let
$\tau(\Breed_i)$ be the per-invocation latency of breed $\Breed_i$.

\begin{definition}[Ensemble Size Function]
\label{def:ensemble-size}
Given a decision budget $T_D > 0$ and a breed $\Breed_i$ with latency
$\tau(\Breed_i) > 0$, the \emph{ensemble size function} is
\[
  k(T_D,\, \tau(\Breed_i)) \;=\; \left\lfloor \frac{T_D}{\tau(\Breed_i)} \right\rfloor.
\]
For a heterogeneous ensemble of breeds $\{\Breed_1, \ldots, \Breed_m\}$ with respective
latencies $\tau_1, \ldots, \tau_m$, the total invocations affordable within $T_D$ is
\[
  K(T_D) \;=\; \sum_{i=1}^{m} k(T_D,\, \tau_i).
\]
\end{definition}

The ensemble size function is the bridge between hardware speed and epistemic capability:
as $\tau(\Breed_i) \to 0$, $k(T_D, \tau_i) \to \infty$, and the system passes through
qualitative capability thresholds that we characterise precisely in
Section~\ref{sec:phase-transition}.

\begin{example}[Classical vs.\ wasm4pm Regimes]
\label{ex:classical-vs-wasm}
The table below contrasts the ensemble size achievable by the classical MYCIN implementation
(1974, Lisp on a DEC-10) against the wasm4pm WASM-compiled breed with a representative
decision budget of $T_D = 60$ seconds.

\begin{center}
\begin{tabular}{lcccc}
\hline
\textbf{System} & \textbf{Breed} & \textbf{$\tau$ (approx.)} & \textbf{$T_D$} & \textbf{$k(T_D,\tau)$} \\
\hline
Classical MYCIN (1974)  & \texttt{production\_rules} & $10$--$60$ min & $60$ s  & $0$--$0$   \\
wasm4pm v26.6.x         & \texttt{production\_rules} & $\approx 2\,\mu\text{s}$ & $60$ s  & $\approx 3 \times 10^7$ \\
wasm4pm v26.6.x         & \texttt{strips\_planning}  & $\approx 10\,\mu\text{s}$ & $60$ s  & $\approx 6 \times 10^6$ \\
wasm4pm v26.6.x         & \texttt{soar\_rule\_match} & $\approx 5\,\mu\text{s}$  & $60$ s  & $\approx 1.2 \times 10^7$ \\
\hline
\end{tabular}
\end{center}

\smallskip
\noindent
\emph{Note on latency values.}  The $\tau$ figures above are approximate and are presented
to convey order-of-magnitude differences between regimes.  Verified measurements obtained
from Criterion micro-benchmarks are recorded in
\texttt{docs/benchmarks/breed-latency-2026-06-10.md}; the corresponding benchmark source
lives in \texttt{crates/wasm4pm-cognition/benches/breed\_latency.rs}.  Readers performing
replication studies should consult those documents rather than the approximate figures in
this table.

Classical MYCIN could not even complete a single consultation within the 60-second budget
when operating interactively; wasm4pm can perform tens of millions of breed invocations
within the same window.  The implication is not merely quantitative: $k = 0$ and
$k = 3 \times 10^7$ support categorically different epistemic operations.
\end{example}

%% -----------------------------------------------------------------------
\section{The Speed Phase Transition Theorem}
\label{sec:phase-transition}
%% -----------------------------------------------------------------------

We formalise the threshold below which qualitatively new capabilities become available.

\begin{definition}[Critical Latency]
\label{def:critical-latency}
Let $\varepsilon > 0$ be an acceptable error tolerance and let $k^*_\varepsilon$ be the
minimum ensemble size required to achieve error below $\varepsilon$ for a given epistemic
task.  The \emph{critical latency} is
\[
  \tau^*_\varepsilon \;=\; \frac{T_D}{k^*_\varepsilon}.
\]
A breed $\Breed_i$ with $\tau(\Breed_i) \leq \tau^*_\varepsilon$ is said to be
\emph{phase-capable} for the task at tolerance $\varepsilon$.
\end{definition}

\begin{theorem}[Speed Phase Transition]
\label{thm:phase-transition}
Let $\Breed_i$ be a breed satisfying the categorical axioms of
Chapter~\ref{chap:breed} and let $T_D$ be a fixed decision budget.  When
$\tau(\Breed_i) \leq \tau^*_\varepsilon$, three capabilities emerge that are logically
inaccessible when $\tau(\Breed_i) > \tau^*_\varepsilon$:

\begin{enumerate}[label=(\roman*)]
  \item \emph{Bayesian sensitivity analysis.}
  \item \emph{Pareto-frontier multi-objective planning.}
  \item \emph{Adversarial robustness auditing.}
\end{enumerate}
\end{theorem}

\begin{proof}
We prove each sub-clause in turn.

\paragraph{(i) Bayesian sensitivity analysis.}
Let $\theta \in \Theta$ be an uncertain parameter of the reasoning context (e.g., the
prior probability that a symptom cluster implies a given diagnosis in the MYCIN breed).
A Bayesian sensitivity analysis asks: across a grid of $k$ values
$\theta_1, \ldots, \theta_k \in \Theta$, how does the posterior distribution
$P(\text{conclusion} \mid \theta_j, \text{evidence})$ vary?

When $\tau(\Breed_i) > T_D / k^*_\varepsilon$, fewer than $k^*_\varepsilon$ evaluations
can be completed within $T_D$; the analyst receives a sparse, non-representative sample
of $\Theta$ and cannot form a reliable estimate of sensitivity.  Formally, the
Berry--Esseen theorem \cite{berry1941,feller1968} bounds the Kolmogorov--Smirnov
distance between the empirical distribution over $k$ samples and the true posterior by
$C / \sqrt{k}$ for a universal constant $C > 0$.  This bound falls below $\varepsilon$ if and only if $k \geq C^2 / \varepsilon^2$,
which requires $\tau(\Breed_i) \leq T_D \varepsilon^2 / C^2 = \tau^*_\varepsilon$.
Thus, reliable Bayesian sensitivity analysis --- one whose error is below $\varepsilon$
with high probability --- is a \emph{threshold phenomenon}: it appears discontinuously as
$\tau$ drops through $\tau^*_\varepsilon$.  Above the threshold the analyst has no
statistically valid sensitivity estimate; below it the full posterior landscape becomes
visible within the decision window.  This discontinuity is the phase transition.

\paragraph{(ii) Pareto-frontier multi-objective planning.}
A STRIPS-family planning problem with $n$ ground actions and $d$ objective dimensions
has a Pareto front of size $O(n^{(d-1)/d})$ in the worst case.  Enumerating the full
Pareto front requires solving at least $n^{(d-1)/d}$ planning instances, each of which
costs $\tau(\Breed_i)$ wall-clock time.  The budget constraint is therefore
$n^{(d-1)/d} \cdot \tau(\Breed_i) \leq T_D$, equivalently
$\tau(\Breed_i) \leq T_D / n^{(d-1)/d}$.

For the logistics test case used in the wasm4pm validation suite (STRIPS logistics,
$n = 100$ ground actions, $d = 3$ objectives), this bound is
$\tau \leq T_D / 100^{2/3} \approx T_D / 21.5$.
Classical STRIPS planners operating at multi-second latencies cannot enumerate the full
front within any realistic decision budget; wasm4pm STRIPS at $\approx 10\,\mu\text{s}$
(see \texttt{docs/benchmarks/breed-latency-2026-06-10.md}) can enumerate the complete
Pareto front for this instance in roughly $0.2$ seconds, well within a 60-second budget.
When $\tau > \tau^*_\varepsilon$, only a strict subset of the Pareto front is reachable;
the planner may commit to a dominated solution without knowing a non-dominated alternative
exists.  The transition from partial to full Pareto coverage is again discontinuous in
$\tau$: it cannot be approached gradually but snaps into place when the latency crosses
the threshold defined by the instance parameters $n$ and $d$.

\paragraph{(iii) Adversarial robustness auditing.}
A robustness audit enumerates perturbations $\delta \in \Delta$ of the input context and
checks whether the breed's conclusion changes under any perturbation.  Let
$|\Delta| = M$ be the size of the perturbation space (e.g., $M = 2^p$ for $p$ binary
feature flips).  An exhaustive audit requires $M$ invocations and therefore costs
$M \cdot \tau(\Breed_i)$ wall-clock time.  For the audit to complete within $T_D$, we
need $\tau(\Breed_i) \leq T_D / M$.

When $\tau > T_D / M$, the audit must sample $\Delta$ uniformly at random, and with
probability at least $1 - (1 - 1/M)^k \approx 1 - e^{-k/M}$ an adversarial perturbation
is missed after $k$ samples.  This probability exceeds $\varepsilon$ only when
$k \geq M \ln(1/\varepsilon)$, requiring $\tau \leq T_D / (M \ln(1/\varepsilon))$.
Below this threshold the audit is exhaustive; above it the audit is necessarily
incomplete and may certify as robust a system that admits a known adversarial input.
The MYCIN CF boundary tests
(\texttt{test\_cf\_threshold\_boundary\_above} and \texttt{test\_cf\_threshold\_boundary\_below}
in \texttt{production\_rules.rs}) implement exactly this perturbation discipline for the
certainty-factor boundary at $\CF = 0.2$, confirming that the breed changes its
conclusion discontinuously at the boundary.  An adversarial auditor operating above
$\tau^*_\varepsilon$ may never sample the neighbourhood of this boundary and may therefore
miss the most dangerous class of adversarial inputs.

This completes the proof that all three capabilities are threshold phenomena governed by
$\tau^*_\varepsilon = T_D / k^*_\varepsilon$.
\end{proof}

%% -----------------------------------------------------------------------
\section{Ensemble Reasoning as a New Epistemic Mode}
\label{sec:ensemble-epistemic-mode}
%% -----------------------------------------------------------------------

Theorem~\ref{thm:phase-transition} has an immediate corollary that reframes the
significance of sub-microsecond breed latency.

\begin{corollary}[Ensemble Reasoning as a New Epistemic Mode]
\label{cor:ensemble-mode}
When all breeds in the wasm4pm kernel satisfy $\tau(\Breed_i) \leq \tau^*_\varepsilon$ for
a shared tolerance $\varepsilon$ and a common decision budget $T_D$, the system operates
in an \emph{ensemble epistemic mode} that is qualitatively distinct from any single-breed
invocation: it simultaneously supports Bayesian sensitivity analysis, full Pareto-frontier
planning, and exhaustive adversarial auditing within every decision cycle.
\end{corollary}

\begin{proof}
By Theorem~\ref{thm:phase-transition}, each of the three capabilities is individually
accessible when $\tau(\Breed_i) \leq \tau^*_\varepsilon$.  Since the breeds are invoked
independently (each WASM call operates on its own linear memory region), the total budget
consumed by running all three epistemic tasks in a single decision cycle is
\[
  T_{\text{total}} \;=\;
    k^*_{\varepsilon,\text{Bayes}} \cdot \tau_{\text{prod}} +
    n^{(d-1)/d} \cdot \tau_{\text{strips}} +
    M \ln(1/\varepsilon) \cdot \tau_{\text{prod}}
  \;\leq\; T_D,
\]
where the final inequality follows from the phase-capability assumption and the fact that
the three task counts are each no greater than $k^*_\varepsilon$ by definition.
Therefore all three tasks complete within $T_D$, and the system simultaneously provides
uncertainty quantification, optimality guarantees, and adversarial certificates --- the
combination that defines the ensemble epistemic mode.
\end{proof}

The ensemble epistemic mode has no classical analogue.  Classical expert systems operated
one breed at a time, produced a single conclusion, and had no mechanism for
uncertainty quantification or adversarial testing within a clinical or operational
decision cycle.  The speed-enabled phase transition is therefore not a performance
improvement; it is a new class of reasoning artifact.

%% -----------------------------------------------------------------------
\section{Summary}
\label{sec:speed-summary}
%% -----------------------------------------------------------------------

This chapter established that execution latency is not merely a performance metric but a
categorical determinant of epistemic capability: below the critical latency $\tau^*_\varepsilon$,
three qualitatively distinct reasoning operations --- Bayesian sensitivity analysis,
Pareto-frontier planning, and adversarial robustness auditing --- become accessible as
phase transitions.  The wasm4pm kernel, whose breed latencies are measured and documented
in \texttt{docs/benchmarks/breed-latency-2026-06-10.md}, operates below this threshold
for all 60 registered breeds, enabling the ensemble epistemic mode of Corollary~\ref{cor:ensemble-mode}.
Chapter~\ref{chap:receipt} turns from capability to accountability: given that the system
now manufactures millions of reasoning conclusions per decision cycle, how is each
conclusion made auditable, non-repudiable, and causally traceable to its inputs?
